Quantização é um método de simplificação dos parâmetros de um modelo de \gls{rna}. 
Como \citeonline{ferreira2022compressao} explica, a aplicação de tal técnica 
traz como benefício a redução do tamanho de modelos, 
viabilizando a utilização em equipamentos de baixo poder computacional, 
além do aumento de velocidade de inferência e simplificação de cálculos realizados pelo processador, 
frente a um modelo que, na maioria dos casos, possui seus parâmetros representados em pontos 
flutuantes de 32 bits. 

A técnica \gls{qat}, com foco em transformação de parâmetros para inteiros de 8 bits,
simula a perda de precisão da conversão de FP32 para INT8 durante o treinamento, 
permitindo que os pesos da rede se adaptem e mitiguem a degradação de acurácia.
Em sistemas de OCR para medidores, 
modelos quantizados alcançam tamanhos da ordem de centenas de kB e 
inferência em subsegundo em hardware embarcado 
\cite{ademola2023resource, jawas2018water}.

